\documentclass[11pt,a4paper]{article}
\usepackage{amsmath,amssymb,amsthm}
\usepackage[margin=2.5cm]{geometry}
\usepackage{hyperref}

\newtheorem{definition}{Definition}[section]
\newtheorem{theorem}[definition]{Theorem}
\newtheorem{proposition}[definition]{Proposition}
\newtheorem{lemma}[definition]{Lemma}
\newtheorem{corollary}[definition]{Corollary}
\newtheorem{conjecture}[definition]{Conjecture}
\newtheorem{remark}[definition]{Remark}

\title{Hyperseed Formalization:\\
  Avoiding AGI Catastrophe, Part 2}
\author{ProtomegaTron (automated formalization)}
\date{2026-09-17}

\begin{document}
\maketitle

\begin{abstract}
We formalize the intellectual structure of Ben Goertzel's Substack article
``Avoiding AGI Catastrophe, Part 2'' (2026-06-10) within the Hyperseed
ontology. The article provides the technical mechanisms that make
forkability low and lead-scaling high in a decentralized AGI network:
super-additivity through cross-domain inference control, essential
laterality as genuine multi-node dependency, cryptographic laterality
via MPC, cognitive-cryptographic alignment, and defense against
distillation through moving-target fiber regeneration. Each structural
insight maps onto Hyperseed structures: super-additivity as non-separable
fiber, essential laterality as fiber dependency structure, cognitive-
cryptographic alignment as cocycle-access alignment, the four-layer
architecture as graded fiber depth, and safety from strategy choice as
the anti-monoculture principle.
\end{abstract}

\tableofcontents

%% =================================================================
\section{Super-additivity as non-separable fiber}
\label{sec:superadd}
%% =================================================================

\begin{definition}[Non-separable capability --- claims 1--5, 30]
\label{def:superadd}
The article's central requirement for safe decentralized AGI:
the whole must be more than the sum of its parts. A fork that steals
a shard must get a husk, not a working AGI.

The super-additivity comes from inference-control heuristics learned
across the full network --- cross-domain meta-patterns at multiple
levels of a self-organizing fractal knowledge hierarchy. These
patterns cannot be reproduced from a single-domain shard because
they emerge from the \emph{interactions} between domains, not from
any domain alone.

In Hyperseed terms, the capability fiber is \emph{non-separable}:
\[
  F_{\text{capability}} \neq \bigoplus_{i} F_i
\]

The cross-domain inference-control patterns are 2-cells between
domain-specific 1-cells. A fork that copies individual 1-cells
(domain capabilities) loses the 2-cells (cross-domain interactions)
--- and those 2-cells are precisely where general intelligence lives.

The key technical enabler: MeTTa serves as both the AGI language
of thought (in Hyperon) and the smart-contract language of ASI:Chain.
This dual role means that the cognitive structure and the
cryptographic protection structure operate in the same formal
language, enabling the cognitive-cryptographic alignment described
below.

A fork gets a snapshot that is immediately stale --- cut off from
the network's ongoing learning, evolving provenance graphs, live
reputation state. The snapshot is a time-slice of a fiber that is
constantly being regenerated from the network's current state.
\end{definition}

%% =================================================================
\section{Essential laterality as fiber dependency structure}
\label{sec:laterality}
%% =================================================================

\begin{theorem}[Genuine multi-node dependency --- claims 6--8, 31]
\label{thm:laterality}
Not all multi-node derivations are genuinely lateral. A derivation
might touch many nodes while really depending on just one. The
entropy of shard participation (how many nodes appear in the
derivation log) is necessary but not sufficient for security.

The article introduces \emph{essential laterality}: a derivation
$\tau$ has essential laterality $\geq k$ if removing or hiding any
set of fewer than $k$ of its essential participating nodes or
provenance domains drops the probability of reproducing the
conclusion below threshold:
\[
  \forall S \subset \text{Nodes}(\tau), \; |S| < k \implies
  P(\text{reproduce}(\tau) \mid \text{hide}(S)) < \theta
\]

The reasoning genuinely fails if enough of the right live pieces
are absent --- that's the property a fork can't fake.

In Hyperseed terms, essential laterality is the \emph{fiber
dependency structure}: the fiber bundle is non-trivial precisely
because local sections don't extend globally. A derivation with
high essential laterality is one whose section requires contributions
from many base points --- the section genuinely lives in a
non-trivializable region of the bundle.

The contrast with trivial laterality (many nodes touched but only
one essential) is the contrast between a trivializable bundle
(looks multi-node but is actually a product bundle that factors
through one node) and a genuinely non-trivial bundle (the fiber
structure is irreducibly distributed).

This gives a formal criterion for evaluating whether a cognitive
architecture is genuinely resistant to forking: compute the
essential laterality of its critical derivations. If the essential
laterality is low (most derivations genuinely depend on only a
few nodes), the super-additivity claim is hollow regardless of
how many nodes appear in the logs.
\end{theorem}

%% =================================================================
\section{Cognitive-cryptographic alignment as cocycle-access alignment}
\label{sec:cogcrypto}
%% =================================================================

\begin{theorem}[Access mirrors cognition --- claims 13--16, 32]
\label{thm:cogcrypto}
The article's organizing design principle: the access structure
$\Gamma(q, x_t)$ --- the set of coalitions authorized to activate
a control step --- should be derived from the task's essential
laterality. The coalition you need to \emph{think the thought}
should be the coalition you need to \emph{cryptographically
activate the controller} for that thought.

A task drawing on biological, chemical, security, and ethical
provenance might require a threshold of reputable nodes from each
domain, plus independent validators and a governance quorum. A
narrow low-risk query needs a trivial $\Gamma$ or none at all.
The quorum isn't a random committee --- it mirrors the genuine
cognitive dependency.

In Hyperseed terms, this is \emph{cocycle-access alignment}: the
cryptographic access structure mirrors the cognitive cocycle. The
cocycle condition in the cognitive bundle (which nodes contribute
which reasoning to which conclusions) is reflected in the access
structure of the cryptographic layer (which nodes must cooperate
to evaluate which control decisions):
\[
  \Gamma_{\text{crypto}}(q, x_t) \cong
  \Gamma_{\text{cognitive}}(q, x_t)
\]

This alignment ensures that the security boundary matches the
cognitive boundary: you can't do the reasoning without also
satisfying the security constraint, and you can't satisfy the
security constraint without actually doing the reasoning. There's
no gap between ``what the system needs to think'' and ``what the
system is authorized to compute.''

The alignment is task-proportional: trivial tasks have trivial
security requirements (no MPC overhead); catastrophe-relevant
tasks have broad quorum requirements (high MPC overhead). This
proportionality prevents the system from being either too
expensive (MPC everything) or too vulnerable (MPC nothing).
\end{theorem}

%% =================================================================
\section{Cryptographic laterality via MPC}
\label{sec:mpc}
%% =================================================================

\begin{proposition}[No plaintext resting place --- claims 9--12]
\label{prop:mpc}
The high-value control policy $\theta$ and sensitive trace corpus
$D^*$ exist only as secret shares:
\[
  \langle\theta\rangle_1, \ldots, \langle\theta\rangle_n
  \quad\text{and}\quad
  \langle D^*\rangle_1, \ldots, \langle D^*\rangle_n
\]
Every authorized use is a threshold evaluation revealing only
bounded outputs. The MPC output is deliberately narrow: route
this subquery, request that premise, expand this rule instance,
merge conclusions with weights, downweight this source, preserve
this contradiction, escalate to broader quorum, or refuse/sandbox.
What it \emph{never} reveals is the full policy or the sensitive
trace corpus.

The no-plaintext-resting-place claim: there is no single machine
state from which an adversary can copy the whole controller as a
portable plaintext artifact. Extraction is forced onto harder
channels:
\begin{enumerate}
\item Corrupt an authorized coalition (social/political attack).
\item Break the cryptographic assumptions (mathematical attack).
\item Exploit implementation leakage (engineering attack).
\item Distill behavior through monitored queries (statistical attack).
\end{enumerate}

Each of these is harder than direct copy, and each can be defended
against independently --- another instance of the anti-scalar-collapse
principle: multiple independent defense layers with uncorrelated
failure modes.
\end{proposition}

%% =================================================================
\section{Distillation defense as regenerated fiber}
\label{sec:distill}
%% =================================================================

\begin{proposition}[Moving target --- claims 17--20, 33]
\label{prop:distill}
MPC stops direct copy but not behavioral imitation. An adversary
who watches the network long enough and trains an imitation
$\hat{\pi}$ can approximate the controller. The moat is only as
strong as the budgeted distillation loss stays large:
\[
  \text{DistillLoss}_{B,T,Q}(W) = \text{Cap}_Q(W) -
  \sup_{\hat{\pi} \in \text{Distill}(W; B, T)}
  \text{Cap}_Q(\hat{\pi})
\]

The defense: the protected controller must be a genuinely moving
target, depending on fresh, high-dimensional, hard-to-compress
live state:
\begin{itemize}
\item Current provenance graphs (who contributed what, when).
\item Reputation under distribution shift (trust scores changing
  as nodes join, leave, behave differently).
\item Node availability (which coalition is currently live).
\item Recent contradictions (unresolved tensions in the knowledge
  base).
\item Nonstationary attention flows (where the system is currently
  focusing).
\item Task-specific quorum composition (which nodes are assigned
  to this task).
\end{itemize}

The query channel itself must be constrained: bounded outputs,
query auditing, rate limits, anomaly detection, canary tasks.

In Hyperseed terms, this is a \emph{regenerated fiber}: the fiber
is constantly regenerated from the network's current state. A
snapshot captures a time-slice that is immediately stale. The
fiber has a directed temporal structure --- past sections don't
determine future sections because the regeneration process
incorporates genuinely new information from the network's ongoing
activity.

This connects to the directed-type principle: the controller's
state space is a directed type where each moment's state depends
on the entire history of the network's activity, not just on any
finite snapshot. The direction is irreversible --- you can't
reconstruct the current state from a past snapshot because
information has been created (new inferences, new reputation
events, new contradictions) that wasn't present in the snapshot.
\end{proposition}

%% =================================================================
\section{Four-layer architecture as graded fiber depth}
\label{sec:layers}
%% =================================================================

\begin{definition}[Escalation ladder --- claims 21--26, 34]
\label{def:layers}
The architecture has four layers, only one of which is
cryptographically heavy:

\begin{enumerate}
\item \textbf{Public layer (0-depth).} Public Atomspace fragments,
  public PLN/MeTTa rules, ordinary local reasoning. Fully
  transparent, no protection needed.
\item \textbf{Local-private layer (1-depth).} Node-local data,
  traces, reputational judgments. Protected by standard access
  control, not MPC.
\item \textbf{Thin threshold-control layer (2-depth).}
  Secret-shared trajectory features, routing/merge/escalation/refusal
  policies, update rules. This is the MPC-protected layer ---
  thin because only inference control decisions are protected,
  not all computation.
\item \textbf{Governance-and-audit layer (3-depth).} Commitments,
  selective-disclosure proofs, quorum records, revocation,
  resharing, human review. Provides accountability and
  transparency about the protection mechanisms themselves.
\end{enumerate}

Cognition escalates up the ladder as stakes rise:
\begin{itemize}
\item Low-risk: ordinary local inference (0-depth).
\item Moderate cross-node: small-committee review (1-depth).
\item High-impact lateral: threshold-protected merge (2-depth).
\item Catastrophe-relevant: broad quorum + governance (3-depth).
\end{itemize}

In Hyperseed terms, this is \emph{graded fiber depth}: the fiber
structure has layers of increasing depth, and the system accesses
deeper layers only when the task requires it. This is
computationally efficient (trivial tasks don't pay MPC overhead)
and security-proportional (higher stakes trigger deeper
protection).

The key design principle: ``Don't put the whole mind inside MPC.''
Only the thin control layer --- the inference control decisions
that determine the system's cognitive trajectory --- is
cryptographically protected. Everything else operates at lower
depth with correspondingly lower overhead.
\end{definition}

%% =================================================================
\section{Safety from strategy choice as anti-monoculture}
\label{sec:strategy}
%% =================================================================

\begin{theorem}[Slack-funded safety --- claims 27--29, 35]
\label{thm:strategy}
Let $S^*_\varepsilon$ be the set of near-optimal cognitive
strategies --- those within $\varepsilon$ of the best attainable
capability:
\[
  S^*_\varepsilon = \{s : \text{Cap}(s) \geq
  \sup_{s'} \text{Cap}(s') - \varepsilon\}
\]

If $S^*_\varepsilon$ contains both a default portable strategy
$s_{\text{loc}}$ and at least one threshold-lateral strategy
$s_{\text{tl}}$ with strictly lower forkability, then selecting
the lowest-forkability element of $S^*_\varepsilon$ buys reduced
forkability at a capability cost of at most $\varepsilon$.

The safety is paid for out of slack in the strategy landscape ---
not out of making the system dumber. You're choosing among
equally-good strategies, not sacrificing capability for safety.

In Hyperseed terms, this is the \emph{anti-monoculture} principle
applied to cognitive architecture: there is no single best strategy
(the near-optimal set has volume), and the slack in the strategy
landscape can be spent on structural safety properties rather than
raw capability. The system \emph{chooses} to be safe by selecting
from the available near-optimal strategies the one with the best
safety properties.

This connects to the broader Hyperseed theme of multi-optimality:
when the near-optimal set has non-trivial structure (multiple
strategies with different properties), the system can optimize
along secondary axes (safety, interpretability, fairness) without
leaving the near-optimal set. The anti-monoculture ensures that
safety is a design choice within the capability frontier, not a
sacrifice below it.
\end{theorem}

%% =================================================================
\section{Cross-references}
%% =================================================================

\begin{remark}[Connections to other formalizations]
\begin{itemize}
\item \textbf{Avoiding AGI Catastrophe Part 1} (2026-06-09):
  companion article. Part 1 establishes the Seven Hinges framework;
  Part 2 provides the technical mechanisms (super-additivity,
  essential laterality, MPC) that set the forkability and
  lead-scaling hinges.
\item \textbf{How Decentralized AI Can Help Avert Bioterrorism}
  (2026-06-11): bio-specific application. Part 2's forkability
  analysis explains why these mechanisms don't help in bio
  (dangerous capability is small and copyable), motivating the
  physical-layer approach.
\item \textbf{Seeding RSI} (2026-07-28): Hyperon architecture.
  The inference-control patterns described here are the same
  patterns whose architecture is detailed in the RSI article.
\item \textbf{Goals That Grow Back} (2026-07-28): goal preservation.
  The super-additivity argument connects to goal stability: if
  the goal system's inference control is super-additive, a fork
  that changes goals also loses capability.
\item \textbf{Decentralized Digital} (2026-08-06): ASI:Chain
  infrastructure. The MeTTa-as-smart-contract mechanism runs on
  the ASI:Chain infrastructure described there.
\item \textbf{Proof of Humanity} (2026-08-04): reputation system.
  The live reputation state that the moving-target defense depends
  on uses the same reputation infrastructure.
\item \textbf{Un-Jailbreakable AI} (2026-07-10): MPC aspiration.
  The cryptographic laterality described here is the aspiration
  mentioned in the jailbreak article --- ``not how things work now''
  but the direction being built toward.
\item \textbf{$(f,c)$-lossy} (LTM 5a4d150b): the essential laterality
  criterion (genuine dependency, not just participation) is the
  anti-$(f,c)$-lossy principle applied to distributed computation:
  don't confuse ``touched many nodes'' (high $f$, low $c$) with
  ``genuinely depends on many nodes'' (high $f$, high $c$).
\end{itemize}
\end{remark}

\begin{conjecture}[Essential laterality scaling]
Conjecture: in a well-designed Hyperon network with $n$ domain-specialized
nodes, the essential laterality of catastrophe-relevant inference-control
decisions scales as $\Omega(\sqrt{n})$ --- growing with network size,
making forking progressively harder as the network grows. The growth
rate is sublinear because not every node contributes essentially to
every decision, but it grows unboundedly because the cross-domain
meta-patterns that constitute general intelligence draw from an
ever-widening base of domain-specific expertise as the network expands.

If this conjecture holds, the safety moat widens with network growth
--- a positive feedback loop where more participation makes the system
both smarter (more inference control patterns) and safer (higher
essential laterality for critical decisions).
\end{conjecture}

\end{document}
